\chapter{1969:Derek Barton and Odd Hassel}

\label{chap:chemistry1969}

The Nobel Prize in Chemistry for the year 1969 was awarded jointly to Derek Barton and Odd Hassel.

\textbf{Motivation:}

for their contributions to the development of the concept of conformation and its application in chemistry

\section{Derek Barton}

\subsection*{Early Life and Education}

Derek Barton was born on 1918-09-08 in Gravesend, England.

\subsection*{Career and Major Research}

Barton studied at Imperial College London, where he received his doctorate and later became a professor. He introduced conformational analysis, showing that the three-dimensional shape of a molecule determines its chemical reactivity, and applied it to the chemistry of steroids and other natural products. He also discovered the photochemical reaction that bears his name.

\subsection*{Nobel Prize-winning Work}

The work for which Derek Barton was awarded the Nobel Prize in Chemistry in 1969 was described by the Nobel Committee as: "for their contributions to the development of the concept of conformation and its application in chemistry".

Building on Hassel's structural studies of cyclohexane, Barton showed that the preferred conformation of a steroid controls how its substituents react, providing a predictive framework for organic stereochemistry.

\subsection*{Significance and Impact}

Conformational analysis became a fundamental tool in organic chemistry, natural product synthesis, and drug design.

\subsection*{Later Career and Honors}

Barton later directed a laboratory of the CNRS at Gif-sur-Yvette in France. He died on 1998-03-16 in College Station, Texas, USA.

\subsection*{Legacy}

The concepts of axial and equatorial substituents and of conformational control of reactivity remain central to organic chemistry.

\subsection*{Modern Developments}

Barton's conformational analysis explained how molecular shape controls reaction rates and selectivity, becoming a cornerstone of modern organic and medicinal chemistry. His ideas guide the rational design of drugs, natural product syntheses, and stereoselective reactions. Barton's later radical chemistry, including the Barton reaction, remains widely used in synthesis.

\subsection*{Selected Publications}

"The Conformation of the Steroid Nucleus" (Experientia, 1950)

"The Principles of Conformational Analysis" (Nobel Lecture, 1969)

\clearpage

\section{Odd Hassel}

\subsection*{Early Life and Education}

Odd Hassel was born on 1897-05-17 in Oslo, Norway.

\subsection*{Career and Major Research}

Hassel studied at the University of Oslo and earned his doctorate in Berlin, then returned to Norway, where he was professor at the University of Oslo. Using electron diffraction and dipole moment measurements, he established that cyclohexane exists preferentially in the chair conformation with distinct axial and equatorial positions.

\subsection*{Nobel Prize-winning Work}

The work for which Odd Hassel was awarded the Nobel Prize in Chemistry in 1969 was described by the Nobel Committee as: "for their contributions to the development of the concept of conformation and its application in chemistry".

Hassel's structural determinations of cyclohexane and its derivatives provided the experimental foundation for Barton's conformational analysis of organic molecules.

\subsection*{Significance and Impact}

Hassel's results showed that saturated six-membered rings adopt well-defined conformations, a fact essential to understanding the structure and reactivity of alicyclic compounds.

\subsection*{Later Career and Honors}

Hassel remained at the University of Oslo for his entire career. He died on 1981-05-11 in Oslo, Norway.

\subsection*{Legacy}

The chair and boat description of cyclohexane introduced through Hassel's work is standard content in every organic chemistry course.

\subsection*{Modern Developments}

Hassel's studies of cyclohexane conformation established the concept of axial and equatorial positions that is now fundamental to understanding the shape and reactivity of organic molecules. Conformational analysis underpins the design of pharmaceuticals, polymers, and molecular machines, and Hassel's insights remain central to organic stereochemistry.

\subsection*{Selected Publications}

"Structural Aspects of Interatomic Charge-Transfer Bonding" (Nobel Lecture, 1969)

\clearpage

\section*{Chapter Summary}

The 1969 prize recognized Barton and Hassel for developing the concept of conformation and applying it in chemistry. Their work linked molecular shape to chemical reactivity and founded conformational analysis.

